\section{Functional Requirements}
\label{sec:functional-requirements}
The catalogue groups related behaviours into \identifier{FR-001} to \identifier{FR-064}.
Each ID groups a related capability; its lettered criteria are the independently
referenced verification conditions. Groups organise the catalogue by affinity,
while shared constraints are stated once.
The catalogue defines the working functional baseline for the single final
delivery in Section~\ref{sec:final-delivery}; stakeholder approval and verification
remain distinct from inclusion.
Requirements retain stable identifiers across revisions. \identifier{FR-063} and \identifier{FR-064}
specify document generation and account updates, respectively.
\subsection{Requirements Catalogue}
\label{sec:requirements-catalogue}
\paragraph{Common acceptance conditions.}
The shared conditions in Section~\ref{sec:shared-definitions} apply to
every relevant requirement: business/operation authorisation, accurate
local/server status, retry without duplicate effects, preservation of
accepted data on cancellation, recoverable validation failures and
explicit handling of unknown values. Connectivity remains operation-specific.
The server revalidates offline proposals, including access and conflicts.
Shared conditions remain mandatory without repetition in every row.
A criterion referencing a business rule or section incorporates its applicable
conditions; verification records the criterion and those conditions together.
\paragraph{Criterion references for verification.}
Verification checks and implementation tasks identify the requirement and
the relevant lettered acceptance criterion. For example, \identifier{FR-041(a)} covers
acceptance of a new current version, \identifier{FR-041(b)} covers retention and deletion,
and \identifier{FR-041(c)} covers incorrect uploads and withdrawal.
\Needspace{10\baselineskip}
\subsubsection{Registration, Authentication and Access}
{\requirementtable
\begin{longtable}{L{0.18}L{0.82}}
\caption{Functional requirements: Registration, Authentication and Access.}\label{tab:catalogue-access}\\
\toprule\rowcolor{TableHeader}
\textbf{Reference} & \textbf{Requirement and acceptance criteria}\\
\midrule\endfirsthead
\caption[]{Functional requirements: Registration, Authentication and Access (continued).}\\
\toprule\rowcolor{TableHeader}
\textbf{Reference} & \textbf{Requirement and acceptance criteria}\\
\midrule\endhead
\midrule\multicolumn{2}{r}{\footnotesize\itshape Continued on next page}\\\endfoot
\bottomrule\endlastfoot
\rowcolor{TableHeader}
\phantomsection\label{fr:001}\textbf{\identifier{FR-001}} & \textbf{The system shall let a Technical Administrator provision the initial owner and private client workspace.}\\*
\phantomsection\label{ac:001:a}\identifier{FR-001(a)} & Require online technical authority, owner name/email and business name/timezone; enforce unique normalised email and one initial owner/workspace, including concurrent retries. Public owner registration is unavailable.\\
\phantomsection\label{ac:001:b}\identifier{FR-001(b)} & Send a single-use 24-hour activation link under \identifier{SEC-09}. The recipient verifies email and chooses their password; controlled resend invalidates the previous link. Never distribute a password in the message.\\
\phantomsection\label{ac:001:c}\identifier{FR-001(c)} & Activate the existing pending account/membership without duplicate setup or operational records. Technical administration grants no business-data rights. Manager accounts continue through valid owner invitations.\\
\midrule
\rowcolor{TableHeader}
\phantomsection\label{fr:002}\textbf{\identifier{FR-002}} & \textbf{The system shall authenticate account holders and enforce business access.}\\*
\phantomsection\label{ac:002:a}\identifier{FR-002(a)} & Only activated accounts may log in with email and password; invalid credentials must not unnecessarily disclose account information.\\
\phantomsection\label{ac:002:b}\identifier{FR-002(b)} & Validate system authority on technical endpoints; validate membership, scope and operation grants on business endpoints. Route administrators to technical account setup and business members to authorised workspaces; apply session expiry.\\
\midrule
\rowcolor{TableHeader}
\phantomsection\label{fr:003}\textbf{\identifier{FR-003}} & \textbf{The system shall terminate sessions, isolate accounts and clean eligible local copies.}\\*
\phantomsection\label{ac:003:a}\identifier{FR-003(a)} & Online or offline logout blocks local access. When connected, identify pending work and offer synchronisation first.\\
\phantomsection\label{ac:003:b}\identifier{FR-003(b)} & Preserve pending work in protected, account-isolated storage; other accounts cannot access or submit it. Explicit discard requires the originating authenticated user's confirmation of the identified losses.\\
\phantomsection\label{ac:003:c}\identifier{FR-003(c)} & Normal recovery requires the originating account and renewed authentication/authorisation. Deferred server-side logout completes on reconnection.\\
\phantomsection\label{ac:003:d}\identifier{FR-003(d)} & Remove eligible centrally accepted local copies after 30 days of account/business/device inactivity. Check on opening before resetting last use; authorised use resets it.\\
\phantomsection\label{ac:003:e}\identifier{FR-003(e)} & Inactivity cleanup never discards unrelated pending work or recovery dependencies. It does not alter the separate current-source replacement policy, central backup expiry or offline-authorisation limit. Replaced-source signing remains ineligible; removed accepted data needs online preparation again.\\
\midrule
\rowcolor{TableHeader}
\phantomsection\label{fr:004}\textbf{\identifier{FR-004}} & \textbf{The system shall enforce bounded offline continuity and revalidate access.}\\*
\phantomsection\label{ac:004:a}\identifier{FR-004(a)} & Prepared offline access lasts 15 days from successful online permission validation; only another successful online validation renews it.\\
\phantomsection\label{ac:004:b}\identifier{FR-004(b)} & Expiry blocks consultation/new operations until connected authentication and validation; preserve pending work and dependencies. Enforce communicated revocation without claiming instant enforcement offline.\\
\phantomsection\label{ac:004:c}\identifier{FR-004(c)} & Reconnection checks uncertain outcomes and revalidates membership, scope, permissions, revisions and dependencies before safe synchronisation. Keep the separate 30-day cache policy in \identifier{FR-003}.\\
\phantomsection\label{ac:004:d}\identifier{FR-004(d)} & Support pending campaign starts and new workers, preserving dependent operations and conflicts without silent overwrite. Present routine status under Section~\ref{sec:offline-interaction}.\\
\midrule
\rowcolor{TableHeader}
\phantomsection\label{fr:005}\textbf{\identifier{FR-005}} & \textbf{The system shall allow users to recover their account access through a verified reset process.}\\*
\phantomsection\label{ac:005:a}\identifier{FR-005(a)} & A recovery request gives a generic acknowledgement; eligible accounts receive a reset message.\\
\phantomsection\label{ac:005:b}\identifier{FR-005(b)} & Only valid, unexpired, unused authorisation permits a policy-compliant reset; consume it and invalidate online sessions on success.\\
\phantomsection\label{ac:005:c}\identifier{FR-005(c)} & Preserve ownership/permissions. Offline devices detect invalidation on reconnection or authorisation expiry; recovery without inbox access is not assumed.\\
\midrule
\rowcolor{TableHeader}
\phantomsection\label{fr:006}\textbf{\identifier{FR-006}} & \textbf{The system shall invite managers and manage invitation acceptance.}\\*
\phantomsection\label{ac:006:a}\identifier{FR-006(a)} & The owner sets recipient email, scope and permissions and may inspect, resend or cancel pending invitations.\\
\phantomsection\label{ac:006:b}\identifier{FR-006(b)} & Only the intended verified account may accept a valid invitation. Invitee registration does not automatically create an owner workspace.\\
\phantomsection\label{ac:006:c}\identifier{FR-006(c)} & Repeated acceptance/resending creates no duplicate membership; expired or cancelled invitations grant no access.\\
\midrule
\rowcolor{TableHeader}
\phantomsection\label{fr:007}\textbf{\identifier{FR-007}} & \textbf{The system shall configure delegated scope and enforce the initial manager profile.}\\*
\phantomsection\label{ac:007:a}\identifier{FR-007(a)} & Only the owner assigns properties and independent operation grants. Scope includes their parcels and current/future campaigns until revoked; no independent campaign/parcel scope grants apply. Managers cannot delegate authority.\\
\phantomsection\label{ac:007:b}\identifier{FR-007(b)} & Restrict multi-property campaign records/indicators to assigned properties. Expose relevant workers/insurance and owner-explicitly available workers, without revealing their other-property work or full mixed-property delivery/file totals outside the grant.\\
\phantomsection\label{ac:007:c}\identifier{FR-007(c)} & Apply the defaults in Section~\ref{sec:permission-defaults}. Financial/full-ticket consultation is disabled initially and requires an explicit owner grant; financial entry/correction remain separate. Evidence upload uses its entry/attachment grant, even when the source contains amounts.\\
\phantomsection\label{ac:007:d}\identifier{FR-007(d)} & Keep financial and complete-ticket consultation consistent as full-content access. Apply revocation to prepared data/files without redaction; enforce saved grants online and on reconnection.\\
\midrule
\rowcolor{TableHeader}
\phantomsection\label{fr:008}\textbf{\identifier{FR-008}} & \textbf{The system shall revoke manager access and review existing pending proposals.}\\*
\phantomsection\label{ac:008:a}\identifier{FR-008(a)} & Revocation blocks normal online business access without deleting accounts or records; offline devices enforce it on reconnection or expiry.\\
\phantomsection\label{ac:008:b}\identifier{FR-008(b)} & After origin-account authentication, submit preserved proposals through restricted recovery without normal business access. The owner accepts or rejects each; acceptance is a new explicit authorisation, subject to date, lifecycle, coverage, dependency and conflict checks.\\
\phantomsection\label{ac:008:c}\identifier{FR-008(c)} & Preserve originator and accepting owner separately. Do not infer prior creation or permission from device timestamps; unresolved proposals remain protected under \identifier{BR-004}. Signing-source eligibility cannot be overridden.\\
\midrule
\rowcolor{TableHeader}
\phantomsection\label{fr:064}\textbf{\identifier{FR-064}} & \textbf{The system shall update verified email and password while preserving account identity.}\\*
\phantomsection\label{ac:064:a}\identifier{FR-064(a)} & Connected authenticated users may change passwords after identity revalidation under the configured password/session policy.\\
\phantomsection\label{ac:064:b}\identifier{FR-064(b)} & Validate/verify the normalised new email before activation; retain the old address until then. Failed/expired links and retries change no ownership, membership or authorship.\\
\phantomsection\label{ac:064:c}\identifier{FR-064(c)} & Preserve account ID, ownership, memberships, signatures and authorship; retries create no replacement identity.\\
\end{longtable}
}

\Needspace{10\baselineskip}
\subsubsection{Properties and Campaigns}
{\requirementtable
\begin{longtable}{L{0.18}L{0.82}}
\caption{Functional requirements: Properties and Campaigns.}\label{tab:catalogue-properties-campaigns}\\
\toprule\rowcolor{TableHeader}
\textbf{Reference} & \textbf{Requirement and acceptance criteria}\\
\midrule\endfirsthead
\caption[]{Functional requirements: Properties and Campaigns (continued).}\\
\toprule\rowcolor{TableHeader}
\textbf{Reference} & \textbf{Requirement and acceptance criteria}\\
\midrule\endhead
\midrule\multicolumn{2}{r}{\footnotesize\itshape Continued on next page}\\\endfoot
\bottomrule\endlastfoot
\rowcolor{TableHeader}
\phantomsection\label{fr:009}\textbf{\identifier{FR-009}} & \textbf{The system shall create properties with required and optional descriptive information.}\\*
\phantomsection\label{ac:009:a}\identifier{FR-009(a)} & Require name and textual location; declared hectares and olive varieties are optional and may be added later.\\
\phantomsection\label{ac:009:b}\identifier{FR-009(b)} & Unknown area stays unknown; derive ownership from the workspace. Use textual location and optional declared hectares; maps and map-derived areas are a possible future enhancement outside the final-delivery commitment under \decisionref{Q-SCOPE-02}.\\
\midrule
\rowcolor{TableHeader}
\phantomsection\label{fr:010}\textbf{\identifier{FR-010}} & \textbf{The system shall create named parcels within a property.}\\*
\phantomsection\label{ac:010:a}\identifier{FR-010(a)} & Require a name and authorised parent property; area and varieties are optional.\\
\phantomsection\label{ac:010:b}\identifier{FR-010(b)} & Parcels are optional. Campaign setup creates reusable property parcels under the same grant, not campaign-only divisions.\\
\midrule
\rowcolor{TableHeader}
\phantomsection\label{fr:011}\textbf{\identifier{FR-011}} & \textbf{The system shall allow authorised users to consult and update properties and parcels.}\\*
\phantomsection\label{ac:011:a}\identifier{FR-011(a)} & Preserve identity and historical associations; new parcels do not redistribute earlier property-level records.\\
\phantomsection\label{ac:011:b}\identifier{FR-011(b)} & Prepared details are consultable offline; creation/changes require connectivity. Missing local data does not establish non-existence.\\
\midrule
\rowcolor{TableHeader}
\phantomsection\label{fr:012}\textbf{\identifier{FR-012}} & \textbf{The system shall archive and reactivate properties and parcels.}\\*
\phantomsection\label{ac:012:a}\identifier{FR-012(a)} & Confirmed archival preserves history and blocks new operational associations through the property and its parcels.\\
\phantomsection\label{ac:012:b}\identifier{FR-012(b)} & Reactivation requires connected authority; parent reactivation does not undo independent parcel archival.\\
\midrule
\rowcolor{TableHeader}
\phantomsection\label{fr:013}\textbf{\identifier{FR-013}} & \textbf{The system shall delete eligible properties and parcels.}\\*
\phantomsection\label{ac:013:a}\identifier{FR-013(a)} & Delete only with connected authority, confirmation and no associated records.\\
\phantomsection\label{ac:013:b}\identifier{FR-013(b)} & Otherwise refuse and offer archival; cancellation preserves the record.\\
\midrule
\rowcolor{TableHeader}
\phantomsection\label{fr:014}\textbf{\identifier{FR-014}} & \textbf{The system shall create Planned campaigns and select their coverage.}\\*
\phantomsection\label{ac:014:a}\identifier{FR-014(a)} & Require name, planned start and eligible property/parcel coverage; allow multiple properties. Create Planned campaigns even where other campaigns are planned or executing.\\
\phantomsection\label{ac:014:b}\identifier{FR-014(b)} & Prevent duplicate parent/child coverage. Selection reserves no execution; apply \identifier{FR-019} at start/reopening. New parcels remain reusable.\\
\midrule
\rowcolor{TableHeader}
\phantomsection\label{fr:015}\textbf{\identifier{FR-015}} & \textbf{The system shall allow authorised users to consult and update campaigns and coverage.}\\*
\phantomsection\label{ac:015:a}\identifier{FR-015(a)} & Prepared details are consultable offline; management changes require connectivity.\\
\phantomsection\label{ac:015:b}\identifier{FR-015(b)} & Validate added coverage. Removal requires reopening if closed and resolving dependencies through permitted correction/deletion; unresolved dependencies block removal, without cascade deletion.\\
\midrule
\rowcolor{TableHeader}
\phantomsection\label{fr:016}\textbf{\identifier{FR-016}} & \textbf{The system shall explicitly start campaigns and maintain distinct planned and actual dates.}\\*
\phantomsection\label{ac:016:a}\identifier{FR-016(a)} & Explicit start confirms a separate actual start and changes Planned to In progress; validate authority, chronology and exclusive coverage.\\
\phantomsection\label{ac:016:b}\identifier{FR-016(b)} & Allow prepared authorised offline start. Save a stable pending operation; dependent work/deliveries await ordered server acceptance.\\
\phantomsection\label{ac:016:c}\identifier{FR-016(c)} & Without an actual start, first actual assignment may propose its working date, but starting remains explicit. Planning, presence and deliveries never start campaigns silently; recording time proves neither actual start nor acceptance.\\
\phantomsection\label{ac:016:d}\identifier{FR-016(d)} & Atomically validate current permissions and competing execution on reconnection. Preserve conflicting starts/fieldwork for resolution; never accept both or silently change actual dates.\\
\midrule
\rowcolor{TableHeader}
\phantomsection\label{fr:017}\textbf{\identifier{FR-017}} & \textbf{The system shall close and reopen campaigns with execution validation.}\\*
\phantomsection\label{ac:017:a}\identifier{FR-017(a)} & Connected authorised closure requires an In progress campaign, actual end and activity validation; retain records and \identifier{FR-035} and \identifier{FR-054} completion exceptions.\\
\phantomsection\label{ac:017:b}\identifier{FR-017(b)} & Reopening returns Closed to In progress only after authority, actual-date and exclusive-coverage checks; it bypasses no correction or execution constraint.\\
\phantomsection\label{ac:017:c}\identifier{FR-017(c)} & Confirmed-work correction/removal requires reopening without a grace period. Planned overlap is allowed; conflicting In progress execution blocks reopening.\\
\midrule
\rowcolor{TableHeader}
\phantomsection\label{fr:018}\textbf{\identifier{FR-018}} & \textbf{The system shall archive, unarchive or delete eligible campaigns.}\\*
\phantomsection\label{ac:018:a}\identifier{FR-018(a)} & Archival changes visibility, not lifecycle or execution coverage; unarchiving does not reopen.\\
\phantomsection\label{ac:018:b}\identifier{FR-018(b)} & Delete only without operational records and with connected authority and confirmation; otherwise offer archival.\\
\midrule
\rowcolor{TableHeader}
\phantomsection\label{fr:019}\textbf{\identifier{FR-019}} & \textbf{The system shall prevent conflicting campaign execution without blocking preparation.}\\*
\phantomsection\label{ac:019:a}\identifier{FR-019(a)} & Allow overlapping Planned campaigns; starting, reopening or extending In progress coverage requires exclusive execution in the affected area.\\
\phantomsection\label{ac:019:b}\identifier{FR-019(b)} & Whole-property execution conflicts with every parcel; distinct parcels may run concurrently. Archival never releases In progress coverage.\\
\phantomsection\label{ac:019:c}\identifier{FR-019(c)} & Validate actual chronology and historical overlap, including retrospective changes. Central validation accepts at most one conflicting execution and preserves rejected proposals.\\
\end{longtable}
}

\Needspace{10\baselineskip}
\subsubsection{Workers, Planning and Actual Work}
{\requirementtable
\begin{longtable}{L{0.18}L{0.82}}
\caption{Functional requirements: Workers, Planning and Actual Work.}\label{tab:catalogue-workforce}\\
\toprule\rowcolor{TableHeader}
\textbf{Reference} & \textbf{Requirement and acceptance criteria}\\
\midrule\endfirsthead
\caption[]{Functional requirements: Workers, Planning and Actual Work (continued).}\\
\toprule\rowcolor{TableHeader}
\textbf{Reference} & \textbf{Requirement and acceptance criteria}\\
\midrule\endhead
\midrule\multicolumn{2}{r}{\footnotesize\itshape Continued on next page}\\\endfoot
\bottomrule\endlastfoot
\rowcolor{TableHeader}
\phantomsection\label{fr:020}\textbf{\identifier{FR-020}} & \textbf{The system shall create identified worker records online or offline and maintain existing profiles.}\\*
\phantomsection\label{ac:020:a}\identifier{FR-020(a)} & Require name, telephone and identification type (DNI, NIE or Other with description), issuing country and number. Preserve the number as text and generate a stable internal ID. Creation provides no account or identity verification and requires no scan.\\
\phantomsection\label{ac:020:b}\identifier{FR-020(b)} & Record machinery capability and optional notes separately from tasks; apply same-business duplicate warnings and reviewed distinct records under \identifier{BR-014}, without cross-business merging.\\
\phantomsection\label{ac:020:c}\identifier{FR-020(c)} & Allow authorised same-day offline creation with required fields and a stable pending ID; permit local attendance, work and driving. Documents require accepted individual source preparation under \identifier{FR-045(a)}.\\
\phantomsection\label{ac:020:d}\identifier{FR-020(d)} & Accept worker identity before dependants; recheck duplicates/authority and preserve work/files during resolution. Retries reuse the same worker and operation IDs.\\
\phantomsection\label{ac:020:e}\identifier{FR-020(e)} & Existing-profile administration remains online, preserving associations, current values and modification time. Required-field checks and duplicate warnings apply; country-specific document-format validation is outside the final delivery under \decisionref{Q-ID-01}.\\
\midrule
\rowcolor{TableHeader}
\phantomsection\label{fr:021}\textbf{\identifier{FR-021}} & \textbf{The system shall inactivate and reactivate workers.}\\*
\phantomsection\label{ac:021:a}\identifier{FR-021(a)} & Preserve work/documents, exclude inactive workers from new-work selection and flag future plans without changing them.\\
\phantomsection\label{ac:021:b}\identifier{FR-021(b)} & Reactivation reuses identity; daily non-selection or absence is not inactivation.\\
\midrule
\rowcolor{TableHeader}
\phantomsection\label{fr:022}\textbf{\identifier{FR-022}} & \textbf{The system shall delete eligible worker records.}\\*
\phantomsection\label{ac:022:a}\identifier{FR-022(a)} & Delete only without associated records and with connected authority and confirmation.\\
\phantomsection\label{ac:022:b}\identifier{FR-022(b)} & Otherwise refuse and offer inactivation; cancellation preserves the worker.\\
\midrule
\rowcolor{TableHeader}
\phantomsection\label{fr:023}\textbf{\identifier{FR-023}} & \textbf{The system shall allow authorised users to consult worker profiles and their actual days and hours worked, including available offline data.}\\*
\phantomsection\label{ac:023:a}\identifier{FR-023(a)} & Consult prepared authorised worker profiles and locally created worker records offline, showing pending/availability limits. Missing local data does not prove central non-existence.\\
\phantomsection\label{ac:023:b}\identifier{FR-023(b)} & Profile access alone grants no document rights. \identifier{FR-020} permits offline creation; existing-profile administration remains online.\\
\phantomsection\label{ac:023:c}\identifier{FR-023(c)} & The profile's activity view lists actual working dates, recorded start/end periods, tasks, areas and counted hours per day and overall. Filter by date range and optionally campaign/property/parcel. Apply Section~\ref{sec:worker-activity-calculations}; count each working date once across matching areas/campaigns.\\
\phantomsection\label{ac:023:d}\identifier{FR-023(d)} & Show presence with incomplete periods separately. Plans establish no actual work; missing hours are not zero. Show local pending changes and accepted values with their actual status.\\
\phantomsection\label{ac:023:e}\identifier{FR-023(e)} & Require work-record consultation authority for the displayed scope. Profile access reveals no other-property work. Offline history uses only prepared authorised records, labelling partial coverage and freshness; unavailable history is not an empty result.\\
\midrule
\rowcolor{TableHeader}
\phantomsection\label{fr:024}\textbf{\identifier{FR-024}} & \textbf{The system shall prepare independent daily work plans with group and individual values.}\\*
\phantomsection\label{ac:024:a}\identifier{FR-024(a)} & Planning is optional. Select campaign, eligible area, workers and dates; create an independent plan for each date.\\
\phantomsection\label{ac:024:b}\identifier{FR-024(b)} & Hours/tasks are optional for one area per worker/date; multiple areas require periods for every relevant assignment.\\
\phantomsection\label{ac:024:c}\identifier{FR-024(c)} & Apply group periods, tasks and optional counted/non-counted breaks with individual overrides. Plans confirm neither attendance nor work.\\
\midrule
\rowcolor{TableHeader}
\phantomsection\label{fr:025}\textbf{\identifier{FR-025}} & \textbf{The system shall reuse a plan or actual work as a reviewed planning proposal.}\\*
\phantomsection\label{ac:025:a}\identifier{FR-025(a)} & Copy workers, areas, periods, tasks and applicable breaks to reviewed target-date proposals; preserve source records and do not copy attendance/actual confirmation.\\
\phantomsection\label{ac:025:b}\identifier{FR-025(b)} & Validate each target without silent overwrite/duplication; review inactive workers for eligible replacement or reactivation.\\
\midrule
\rowcolor{TableHeader}
\phantomsection\label{fr:026}\textbf{\identifier{FR-026}} & \textbf{The system shall validate planned periods and scheduling conflicts.}\\*
\phantomsection\label{ac:026:a}\identifier{FR-026(a)} & Defined periods must end after their start; reject same-worker planned overlaps across properties, campaigns and adjacent dates.\\
\phantomsection\label{ac:026:b}\identifier{FR-026(b)} & Allow equal consecutive boundaries; use complete overnight timestamps and start-date association. Missing hours cannot establish conflict-free scheduling.\\
\phantomsection\label{ac:026:c}\identifier{FR-026(c)} & Validate planned and actual interval sets separately; a matching plan does not duplicate actual work.\\
\midrule
\rowcolor{TableHeader}
\phantomsection\label{fr:027}\textbf{\identifier{FR-027}} & \textbf{The system shall adjust eligible plans while preserving their role as an expectation.}\\*
\phantomsection\label{ac:027:a}\identifier{FR-027(a)} & Allow worker removal/replacement before actual work exists in the affected daily-area context.\\
\phantomsection\label{ac:027:b}\identifier{FR-027(b)} & Thereafter preserve that plan as a reference; record differences as actual work/corrections. Other target-date plans remain independently editable.\\
\midrule
\rowcolor{TableHeader}
\phantomsection\label{fr:028}\textbf{\identifier{FR-028}} & \textbf{The system shall allow authorised users to consult schedules through daily, worker and weekly views, including offline preparation.}\\*
\phantomsection\label{ac:028:a}\identifier{FR-028(a)} & Use the same plans in all views with campaign/property/parcel/worker/task filters; identify missing hours/tasks, breaks, inactive workers and overnight ends.\\
\phantomsection\label{ac:028:b}\identifier{FR-028(b)} & Prepare the current and following week; prepare extra dates explicitly. Show freshness and pending changes.\\
\phantomsection\label{ac:028:c}\identifier{FR-028(c)} & Distinguish empty results, failures and unavailable local data; consultation confirms no attendance or work.\\
\midrule
\rowcolor{TableHeader}
\phantomsection\label{fr:029}\textbf{\identifier{FR-029}} & \textbf{The system shall create and adjust plans offline using prepared reference data.}\\*
\phantomsection\label{ac:029:a}\identifier{FR-029(a)} & Save supported plan changes locally as pending; revalidate authority and business rules on reconnection.\\
\phantomsection\label{ac:029:b}\identifier{FR-029(b)} & Rejected/conflicting changes remain unaccepted; preserve source and target-plan integrity during recovery.\\
\midrule
\rowcolor{TableHeader}
\phantomsection\label{fr:030}\textbf{\identifier{FR-030}} & \textbf{The system shall distinguish daily attendance from campaign participation and unrelated workers.}\\*
\phantomsection\label{ac:030:a}\identifier{FR-030(a)} & Keep one attendance decision per worker/start date/business. Selected/expected workers await explicit Present/Absent confirmation; unrelated unselected workers generate neither absence nor alerts.\\
\phantomsection\label{ac:030:b}\identifier{FR-030(b)} & Permit direct actual entry without planning, including newly created offline workers. Confirm campaign-area participation separately; presence may precede hours.\\
\phantomsection\label{ac:030:c}\identifier{FR-030(c)} & Absence cannot coexist with actual assignments/periods. Areas share daily presence; missing-hours alerts require presence and incomplete campaign work. Driver selection creates no presence.\\
\midrule
\rowcolor{TableHeader}
\phantomsection\label{fr:031}\textbf{\identifier{FR-031}} & \textbf{The system shall explicitly confirm actual periods with their task and area.}\\*
\phantomsection\label{ac:031:a}\identifier{FR-031(a)} & Review a plan or enter actual work directly; permit individual changes in group entry without silently confirming attendance or rewriting plans.\\
\phantomsection\label{ac:031:b}\identifier{FR-031(b)} & A confirmed complete period requires task, area and valid start/end instants. Split task/area changes; missing details remain incomplete assignments, not unclassified counted hours.\\
\phantomsection\label{ac:031:c}\identifier{FR-031(c)} & Use the local start date for overnight work. Keep presence, assignment, complete period and synchronisation distinct; first work may explicitly propose start under \identifier{FR-016}.\\
\midrule
\rowcolor{TableHeader}
\phantomsection\label{fr:032}\textbf{\identifier{FR-032}} & \textbf{The system shall record breaks and calculate counted working time.}\\*
\phantomsection\label{ac:032:a}\identifier{FR-032(a)} & Require break start/end with end later than start; record its worker and counted/non-counted treatment. Reject positive-duration overlap with another break for that worker; consecutive boundaries are allowed.\\
\phantomsection\label{ac:032:b}\identifier{FR-032(b)} & Warn when a break lies partly/wholly outside recorded work or hours are still incomplete. Deduct only the non-counted intersection with that worker's actual periods, never counted breaks or already excluded gaps; attribute deductions to affected tasks.\\
\phantomsection\label{ac:032:c}\identifier{FR-032(c)} & Apply Section~\ref{sec:counted-working-time}; no payroll calculation is implied.\\
\midrule
\rowcolor{TableHeader}
\phantomsection\label{fr:033}\textbf{\identifier{FR-033}} & \textbf{The system shall validate actual periods and prevent worker-time overlap.}\\*
\phantomsection\label{ac:033:a}\identifier{FR-033(a)} & Actual periods must end after their start; reject same-worker overlaps across areas, campaigns and adjacent dates.\\
\phantomsection\label{ac:033:b}\identifier{FR-033(b)} & Allow equal consecutive boundaries; reject zero duration and any positive overlap.\\
\phantomsection\label{ac:033:c}\identifier{FR-033(c)} & Use complete overnight timestamps; validate plans separately without counting them as actual work.\\
\midrule
\rowcolor{TableHeader}
\phantomsection\label{fr:034}\textbf{\identifier{FR-034}} & \textbf{The system shall record attendance and actual work offline with its dependencies.}\\*
\phantomsection\label{ac:034:a}\identifier{FR-034(a)} & Use prepared campaign/area context and prepared workers or workers created under \identifier{FR-020}; support pending \identifier{FR-016} starts and retain parent/dependency IDs.\\
\phantomsection\label{ac:034:b}\identifier{FR-034(b)} & Apply the same task, chronology, attendance and overlap checks locally; label provisional totals and revalidate centrally. Preserve work at expiry under \identifier{FR-004}.\\
\midrule
\rowcolor{TableHeader}
\phantomsection\label{fr:035}\textbf{\identifier{FR-035}} & \textbf{The system shall allow authorised users to enter retrospective work and complete missing information.}\\*
\phantomsection\label{ac:035:a}\identifier{FR-035(a)} & Preserve actual working date separately from entry timestamp.\\
\phantomsection\label{ac:035:b}\identifier{FR-035(b)} & After closure, complete missing details/periods within actual dates without changing confirmed values/associations; validate attendance, coverage and overlap under \identifier{BR-011}.\\
\phantomsection\label{ac:035:c}\identifier{FR-035(c)} & Confirmed-work changes/removals require reopening; work outside actual campaign dates also requires date review. Entry never silently changes campaign dates.\\
\midrule
\rowcolor{TableHeader}
\phantomsection\label{fr:036}\textbf{\identifier{FR-036}} & \textbf{The system shall correct attendance, actual work and valid record associations.}\\*
\phantomsection\label{ac:036:a}\identifier{FR-036(a)} & Require correction authority and valid dates, coverage, overlaps and open campaign state in every affected context; no age-based deadline applies.\\
\phantomsection\label{ac:036:b}\identifier{FR-036(b)} & Allow worker, date, campaign, area, period, task and break corrections without changing plans or unrelated work.\\
\phantomsection\label{ac:036:c}\identifier{FR-036(c)} & Changing presence to absence requires confirmed removal of all contradictory assignments, periods and breaks for that business date, including assignments without hours; cancellation preserves them.\\
\phantomsection\label{ac:036:d}\identifier{FR-036(d)} & Moving a period preserves other work and attendance consistency; retain current values without permanent field-change history.\\
\midrule
\rowcolor{TableHeader}
\phantomsection\label{fr:037}\textbf{\identifier{FR-037}} & \textbf{The system shall support offline corrections and authorised conflict resolution.}\\*
\phantomsection\label{ac:037:a}\identifier{FR-037(a)} & Save prepared corrections as pending; compare changed accepted values with the proposal for authorised manual review.\\
\phantomsection\label{ac:037:b}\identifier{FR-037(b)} & Keep accepted values or resubmit a reviewed correction for validation; temporary comparison creates no permanent correction history.\\
\end{longtable}
}

\Needspace{10\baselineskip}
\subsubsection{Worker Documents and Signing}
{\requirementtable
\begin{longtable}{L{0.18}L{0.82}}
\caption{Functional requirements: Worker Documents and Signing.}\label{tab:catalogue-documents}\\
\toprule\rowcolor{TableHeader}
\textbf{Reference} & \textbf{Requirement and acceptance criteria}\\
\midrule\endfirsthead
\caption[]{Functional requirements: Worker Documents and Signing (continued).}\\
\toprule\rowcolor{TableHeader}
\textbf{Reference} & \textbf{Requirement and acceptance criteria}\\
\midrule\endhead
\midrule\multicolumn{2}{r}{\footnotesize\itshape Continued on next page}\\\endfoot
\bottomrule\endlastfoot
\rowcolor{TableHeader}
\phantomsection\label{fr:038}\textbf{\identifier{FR-038}} & \textbf{The system shall allow authorised users to consult and download current authorised worker documents.}\\*
\phantomsection\label{ac:038:a}\identifier{FR-038(a)} & List current contract/insurance/other documents with subtype, title, source identity and individual signing status; show any form response separately and open/download available files. Replaced sources/signatures/PDFs are unavailable; no history browser is provided.\\
\phantomsection\label{ac:038:b}\identifier{FR-038(b)} & Apply business-wide owner access and manager worker/document grants under \identifier{BR-024}; profile access alone is insufficient.\\
\phantomsection\label{ac:038:c}\identifier{FR-038(c)} & Apply \identifier{BR-026} retention and \identifier{BR-025} expiry, including unknown, expired and withdrawn states. Current does not imply valid/signed; unavailable retrieval is not an empty list.\\
\midrule
\rowcolor{TableHeader}
\phantomsection\label{fr:039}\textbf{\identifier{FR-039}} & \textbf{The system shall prepare insurance documents for field consultation.}\\*
\phantomsection\label{ac:039:a}\identifier{FR-039(a)} & Prepare current insurance for workers planned for the current/following week and additional existing workers selected online. Prepare no superseded versions; apply replacement deletion on synchronisation.\\
\phantomsection\label{ac:039:b}\identifier{FR-039(b)} & Show preparation failures, missing files and freshness; apply received access/deletion changes. Preparation coverage does not determine offline expiry.\\
\midrule
\rowcolor{TableHeader}
\phantomsection\label{fr:040}\textbf{\identifier{FR-040}} & \textbf{The system shall upload worker documents and their metadata.}\\*
\phantomsection\label{ac:040:a}\identifier{FR-040(a)} & Online upload requires modification authority, exactly one same-business worker, category, title and an individual file; retain subtype and original contents. Notes are optional; version/upload time are automatic.\\
\phantomsection\label{ac:040:b}\identifier{FR-040(b)} & Apply \identifier{BR-044} formats/approved size limit, subject to representative-file validation; unknown insurance expiry remains unconfirmed.\\
\phantomsection\label{ac:040:c}\identifier{FR-040(c)} & Bind the source to that individual worker. External-signature declarations follow \identifier{BR-027} and create no application signing event.\\
\phantomsection\label{ac:040:d}\identifier{FR-040(d)} & Validate/review before acceptance; incomplete files cannot become current, and retries create no duplicate versions.\\
\midrule
\rowcolor{TableHeader}
\phantomsection\label{fr:041}\textbf{\identifier{FR-041}} & \textbf{The system shall replace current documents, delete superseded files/signatures and manage withdrawal or eligible deletion.}\\*
\phantomsection\label{ac:041:a}\identifier{FR-041(a)} & Keep the current source until a complete replacement is validated and accepted. Failure/cancellation preserves it; retries publish no duplicate replacement. The new source inherits no signature or external declaration.\\
\phantomsection\label{ac:041:b}\identifier{FR-041(b)} & Apply \identifier{BR-026} to all categories: delete the superseded source, signature and signed PDF, including signed contracts and externally declared signed uploads. Retain no consultable version history. Invalidate pending signing of the removed source without transferring or resurrecting it.\\
\phantomsection\label{ac:041:c}\identifier{FR-041(c)} & An authorised user may delete a dependency-free incorrect unsigned upload after confirmation. Signed errors may be withdrawn until replacement; withdrawal does not expose an older version. Successful replacement deletes the withdrawn source and its signing files too.\\
\phantomsection\label{ac:041:d}\identifier{FR-041(d)} & Management is online. Propagate accepted replacement/deletion to prepared copies on synchronisation; keep backup/restore suppression and temporary rejected-proposal recovery separate from normal document history.\\
\midrule
\rowcolor{TableHeader}
\phantomsection\label{fr:042}\textbf{\identifier{FR-042}} & \textbf{The system shall conduct in-person document review and explicit worker signing.}\\*
\phantomsection\label{ac:042:a}\identifier{FR-042(a)} & An authorised initiator identifies the present worker and eligible version; no worker account or owner signature is required.\\
\phantomsection\label{ac:042:b}\identifier{FR-042(b)} & The worker reviews the version, explicitly chooses acceptance/refusal where required, draws and reviews a non-empty signature, and confirms. Redraw/cancel is available before confirmation; cancellation is distinct from signed refusal.\\
\phantomsection\label{ac:042:c}\identifier{FR-042(c)} & Record worker, source version, initiator, time and any required response as immutable evidence while the source is current. A signed refusal is not consent. A new source needs fresh review/signing; successful replacement deletes the old evidence/PDF.\\
\phantomsection\label{ac:042:d}\identifier{FR-042(d)} & Retain source and timestamped worker-specific PDF with signature and any explicit choice in the appropriate form area. Preserve source markings; offline output remains pending until acceptance.\\
\midrule
\rowcolor{TableHeader}
\phantomsection\label{fr:043}\textbf{\identifier{FR-043}} & \textbf{The system shall prevent duplicate signing and recover uncertain submissions.}\\*
\phantomsection\label{ac:043:a}\identifier{FR-043(a)} & Do not repeat completed worker/version signing; show accepted results and check/resume pending or uncertain operations before requesting another signature.\\
\phantomsection\label{ac:043:b}\identifier{FR-043(b)} & Distinguish no confirmed signature, local pending, accepted and resolution-required outcomes; preserve recoverable results without duplicate events.\\
\midrule
\rowcolor{TableHeader}
\phantomsection\label{fr:044}\textbf{\identifier{FR-044}} & \textbf{The system shall bind each document and its signing result to exactly one worker.}\\*
\phantomsection\label{ac:044:a}\identifier{FR-044(a)} & Require exactly one same-business worker per Document. Accepted ownership is immutable; every version, signature and rendition must reference that same worker.\\
\phantomsection\label{ac:044:b}\identifier{FR-044(b)} & Reject another worker signing or receiving a rendition of the document. Generating for several workers creates separate individual documents from the selected template.\\
\midrule
\rowcolor{TableHeader}
\phantomsection\label{fr:045}\textbf{\identifier{FR-045}} & \textbf{The system shall prepare and synchronise offline signing for the reviewed version.}\\*
\phantomsection\label{ac:045:a}\identifier{FR-045(a)} & Prepare and sign authorised individual versions, including contracts, without bulk contract download. A newly created offline worker must first be accepted online and receive an uploaded/generated individual source before offline signing can be prepared.\\
\phantomsection\label{ac:045:b}\identifier{FR-045(b)} & Accept pending signing only against the same current eligible source. Replacement deletes the old source; reject its pending signature and require fresh worker review/signing of the new source.\\
\phantomsection\label{ac:045:c}\identifier{FR-045(c)} & Preserve protected temporary recovery/outcome information for an ineligible source without a version-history browser. Owner approval cannot validate a withdrawn/deleted/replaced source. Never transfer a signature or recreate deleted content; remove superseded prepared copies on synchronisation.\\
\midrule
\rowcolor{TableHeader}
\phantomsection\label{fr:063}\textbf{\identifier{FR-063}} & \textbf{The system shall generate reviewed worker documents from available templates.}\\*
\phantomsection\label{ac:063:a}\identifier{FR-063(a)} & Generate online using the four supplied Spanish templates and their reviewed revision. Prefill known identity and review date, role and form selections in Table~\ref{tab:source-forms}; preserve the Spanish text, identify gaps and leave worker profiles unchanged.\\
\phantomsection\label{ac:063:b}\identifier{FR-063(b)} & Review each document and its worker before unsigned acceptance. Generate a separate individual source per selected worker for every template; preserve accepted versions on failure/cancellation.\\
\phantomsection\label{ac:063:c}\identifier{FR-063(c)} & Use \identifier{UC-14} signing and \identifier{FR-042(d)} PDFs, including prepared offline signing. Required choices are worker-confirmed, never prefilled as consent. Worker days/hours are consulted through \identifier{FR-023}; no dedicated attendance-document generation is included.\\
\end{longtable}
}

\Needspace{12\baselineskip}
\subsubsection{Deliveries, Weighing, Analyses and Reported Values}
{\requirementtable
\begin{longtable}{L{0.18}L{0.82}}
\caption{Functional requirements: Deliveries, Weighing, Analyses and Reported Values.}\label{tab:catalogue-deliveries}\\
\toprule\rowcolor{TableHeader}
\textbf{Reference} & \textbf{Requirement and acceptance criteria}\\
\midrule\endfirsthead
\caption[]{Functional requirements: Deliveries, Weighing, Analyses and Reported Values (continued).}\\
\toprule\rowcolor{TableHeader}
\textbf{Reference} & \textbf{Requirement and acceptance criteria}\\
\midrule\endhead
\midrule\multicolumn{2}{r}{\footnotesize\itshape Continued on next page}\\\endfoot
\bottomrule\endlastfoot
\rowcolor{TableHeader}
\phantomsection\label{fr:046}\textbf{\identifier{FR-046}} & \textbf{The system shall register, correct or cancel deliveries and their provenance, and delete delivery records when permitted.}\\*
\phantomsection\label{ac:046:a}\identifier{FR-046(a)} & Require campaign, eligible origins, registered worker driver, vehicle, destination and departure. Record known source quantities separately from declared transport origin/reference; infer no allocation or driver attendance/hours.\\
\phantomsection\label{ac:046:b}\identifier{FR-046(b)} & Delivery changes require separate correction authority and valid chronology/contexts. Reopen closed affected campaigns for corrections, cancellation, deletion or reassignment; \identifier{FR-054} still covers source-result updates.\\
\phantomsection\label{ac:046:c}\identifier{FR-046(c)} & Confirmed deletion requires no ticket, analysis, valuation or evidence; remove composed transport/batches atomically. Alternatively cancel with a reason, preserving dependencies and excluding indicator contributions under \identifier{BR-031}.\\
\phantomsection\label{ac:046:d}\identifier{FR-046(d)} & Reassignment validates all affected permissions, dates, coverage and dependencies without cascade deletion. Cancel interrupted/refused trips with a reason; preserve their records and exclude contributions. A new physical attempt creates a new delivery/transport record; cancellation is not reversed.\\
\midrule
\rowcolor{TableHeader}
\phantomsection\label{fr:047}\textbf{\identifier{FR-047}} & \textbf{The system shall confirm a delivery or retrospectively register the complete transport.}\\*
\phantomsection\label{ac:047:a}\identifier{FR-047(a)} & Save departure as In transport; authorised arrival confirmation changes it to Delivered with arrival not before departure.\\
\phantomsection\label{ac:047:b}\identifier{FR-047(b)} & Allow retrospective entry as Delivered with both event times, distinct from entry times; one associated transport per delivery remains the proposal.\\
\phantomsection\label{ac:047:c}\identifier{FR-047(c)} & Arrival enables eligible counts/results under \identifier{BR-038}. Offline confirmation updates labelled local provisional indicators; central totals await acceptance.\\
\midrule
\rowcolor{TableHeader}
\phantomsection\label{fr:048}\textbf{\identifier{FR-048}} & \textbf{The system shall record departures and delivery confirmations offline.}\\*
\phantomsection\label{ac:048:a}\identifier{FR-048(a)} & Use prepared campaign, area and worker data; allow a new destination name.\\
\phantomsection\label{ac:048:b}\identifier{FR-048(b)} & Keep physical and synchronisation states separate; validate current authority and campaign context centrally.\\
\midrule
\rowcolor{TableHeader}
\phantomsection\label{fr:049}\textbf{\identifier{FR-049}} & \textbf{The system shall record one weighing ticket per delivery and its reported source values.}\\*
\phantomsection\label{ac:049:a}\identifier{FR-049(a)} & Allow one ticket per non-cancelled delivery, including In transport. Preserve reference, cooperative, source date and Table~\ref{tab:source-ticket} metadata. KILOS ASIGNADOS is the canonical accepted assigned kg for indicators and yield weighting; keep total kg separate, without another editable assigned-weight copy.\\
\phantomsection\label{ac:049:b}\identifier{FR-049(b)} & Record processing start/end and printed duration in seconds; preserve and flag discrepancies. MEDIANERIA is the accepted percentage, not yield. Record any reported excess/penalty with unit and application status; no automatic tariff calculation or repeated percentage/deduction applies.\\
\phantomsection\label{ac:049:c}\identifier{FR-049(c)} & Expose genuinely missing or invalid source values and exclude dependent calculations. Do not infer missing assigned kg from total kg or MEDIANERIA. References are contextual; concurrent creation accepts at most one ticket per delivery.\\
\phantomsection\label{ac:049:d}\identifier{FR-049(d)} & Apply \identifier{BR-044} evidence limits, Section~\ref{sec:permission-defaults} grants and Section~\ref{sec:conceptual-files} transfer recovery. Full-ticket access includes printed amounts without redaction/classification; structured financial writes require separate grants. Prepared authorised data/files work offline.\\
\midrule
\rowcolor{TableHeader}
\phantomsection\label{fr:050}\textbf{\identifier{FR-050}} & \textbf{The system shall correct the current ticket and adjustment information.}\\*
\phantomsection\label{ac:050:a}\identifier{FR-050(a)} & Separate correction authority updates the existing ticket, including post-confirmation completion; retain current values without a second ticket or permanent field history.\\
\phantomsection\label{ac:050:b}\identifier{FR-050(b)} & Refresh indicators under \identifier{FR-061}; preserve uniqueness and use the common conflict workflow.\\
\midrule
\rowcolor{TableHeader}
\phantomsection\label{fr:051}\textbf{\identifier{FR-051}} & \textbf{The system shall record one analysis per delivery and its reported yield.}\\*
\phantomsection\label{ac:051:a}\identifier{FR-051(a)} & Allow one analysis per non-cancelled delivery, including In transport. Record available yield percentage, entity, reference and analysis date separately from entry time; pre-arrival results await receipt under \identifier{BR-038}.\\
\phantomsection\label{ac:051:b}\identifier{FR-051(b)} & Weight yield by the delivery ticket's accepted assigned weight in kg, never an editable analysis copy; exclude missing yield/weight and ineligible or non-positive weight, disclosing incomplete coverage.\\
\phantomsection\label{ac:051:c}\identifier{FR-051(c)} & Allow prepared offline entry and \identifier{BR-044} evidence; track record/file states separately. No automatic extraction occurs.\\
\midrule
\rowcolor{TableHeader}
\phantomsection\label{fr:052}\textbf{\identifier{FR-052}} & \textbf{The system shall correct the current analysis without creating another result for the delivery.}\\*
\phantomsection\label{ac:052:a}\identifier{FR-052(a)} & Separate correction authority updates the existing analysis, including post-confirmation completion; retain one current result without permanent field history.\\
\phantomsection\label{ac:052:b}\identifier{FR-052(b)} & Refresh indicators without changing ticket weights or delivery dates; ticket correction updates the derived reference weight.\\
\midrule
\rowcolor{TableHeader}
\phantomsection\label{fr:053}\textbf{\identifier{FR-053}} & \textbf{The system shall allow authorised users to consult and manage the current reported ticket amount online and offline.}\\*
\phantomsection\label{ac:053:a}\identifier{FR-053(a)} & Require operation-specific financial authority and an existing ticket/non-cancelled delivery, including In transport. Store the EUR amount to two decimal places, the cooperative/date and provisional/final status; reference, notes and PDF/JPG/PNG evidence are optional. Apply \identifier{BR-038} pre-arrival exclusion.\\
\phantomsection\label{ac:053:b}\identifier{FR-053(b)} & Final replaces provisional. A final amount may return to provisional with financial correction permission and explicit confirmation. Retain one current value; missing differs from zero, and changes are never additive monetary contributions.\\
\phantomsection\label{ac:053:c}\identifier{FR-053(c)} & Do not calculate prices or payments, reapply reported penalties or infer area allocations. Source files retain separate transfer/retention status.\\
\phantomsection\label{ac:053:d}\identifier{FR-053(d)} & Allow prepared authorised offline consultation/entry/correction. Apply Section~\ref{sec:permission-defaults} financial/full-file grants; label freshness/pending proposals in indicators/PDFs. No monetary-content classification is required.\\
\midrule
\rowcolor{TableHeader}
\phantomsection\label{fr:054}\textbf{\identifier{FR-054}} & \textbf{The system shall complete and correct external information after campaign closure.}\\*
\phantomsection\label{ac:054:a}\identifier{FR-054(a)} & Authorised ticket, analysis and reported-value additions/corrections for existing deliveries/tickets need no reopening.\\
\phantomsection\label{ac:054:b}\identifier{FR-054(b)} & This does not permit confirmed-work correction, new activity outside actual dates, delivery reversal or campaign reassignment.\\
\end{longtable}
}

\Needspace{14\baselineskip}
\subsubsection{Campaign Indicators and Reports}
{\requirementtable
\begin{longtable}{L{0.18}L{0.82}}
\caption{Functional requirements: Campaign Indicators and Reports.}\label{tab:catalogue-indicators}\\
\toprule\rowcolor{TableHeader}
\textbf{Reference} & \textbf{Requirement and acceptance criteria}\\
\midrule\endfirsthead
\caption[]{Functional requirements: Campaign Indicators and Reports (continued).}\\
\toprule\rowcolor{TableHeader}
\textbf{Reference} & \textbf{Requirement and acceptance criteria}\\
\midrule\endhead
\midrule\multicolumn{2}{r}{\footnotesize\itshape Continued on next page}\\\endfoot
\bottomrule\endlastfoot
\rowcolor{TableHeader}
\phantomsection\label{fr:055}\textbf{\identifier{FR-055}} & \textbf{The system shall allow authorised users to consult campaign indicators and property/parcel breakdowns.}\\*
\phantomsection\label{ac:055:a}\identifier{FR-055(a)} & Consult one campaign during/after activity, with overview, property summaries and available parcel detail.\\
\phantomsection\label{ac:055:b}\identifier{FR-055(b)} & Allow campaigns without activity; preserve identity and label partial authorised scope. Review one campaign at a time. Cross-campaign and cross-year comparisons remain a possible future enhancement outside the final-delivery commitment under \decisionref{Q-SCOPE-03}.\\
\midrule
\rowcolor{TableHeader}
\phantomsection\label{fr:056}\textbf{\identifier{FR-056}} & \textbf{The system shall filter indicators and inspect their supporting records.}\\*
\phantomsection\label{ac:056:a}\identifier{FR-056(a)} & Filter by date/property/parcel/worker/task only where source relationships exist; label scope/date basis. Work uses the business-local period-start date; delivery, weight, yield and reported-value indicators use arrival date. Preserve source/report dates in detail.\\
\phantomsection\label{ac:056:b}\identifier{FR-056(b)} & Enforce permission/local availability for contributing records. Worker/task filters allocate no deliveries/money; unavailable detail is not empty.\\
\midrule
\rowcolor{TableHeader}
\phantomsection\label{fr:057}\textbf{\identifier{FR-057}} & \textbf{The system shall present actual participants, working days and counted hours.}\\*
\phantomsection\label{ac:057:a}\identifier{FR-057(a)} & Count distinct present workers with actual campaign assignment, including incomplete hours, and distinct eligible period-start dates.\\
\phantomsection\label{ac:057:b}\identifier{FR-057(b)} & Show counted hours overall/by worker/task using confirmed complete actual periods and worker breaks, never plans or machinery capability.\\
\phantomsection\label{ac:057:c}\identifier{FR-057(c)} & Recalculate distinct workers/dates across areas; disclose unknown periods and incomplete assignments.\\
\midrule
\rowcolor{TableHeader}
\phantomsection\label{fr:058}\textbf{\identifier{FR-058}} & \textbf{The system shall present delivery, weight and assigned-weighted yield indicators.}\\*
\phantomsection\label{ac:058:a}\identifier{FR-058(a)} & Only Delivered, non-cancelled deliveries contribute. Count them without tickets/analyses; sum each eligible compatible reported weight once.\\
\phantomsection\label{ac:058:b}\identifier{FR-058(b)} & For that set, weight known yield by eligible positive ticket accepted assigned kg. Disclose exclusions, contributing count and incomplete coverage; no eligible denominator means unavailable, not zero.\\
\phantomsection\label{ac:058:c}\identifier{FR-058(c)} & Count mixed loads once; allocate only known source quantities. Under \identifier{BR-035}, mixed-load money remains unallocated by area and percentage allocation is excluded.\\
\phantomsection\label{ac:058:d}\identifier{FR-058(d)} & Show pre-arrival ticket/analysis data as awaiting receipt; exclude it from indicators without falsely labelling it missing.\\
\midrule
\rowcolor{TableHeader}
\phantomsection\label{fr:059}\textbf{\identifier{FR-059}} & \textbf{The system shall present current reported monetary indicators with financial authorisation.}\\*
\phantomsection\label{ac:059:a}\identifier{FR-059(a)} & For Delivered, non-cancelled deliveries, show current provisional/final totals separately and tickets awaiting amounts, online or from authorised prepared data.\\
\phantomsection\label{ac:059:b}\identifier{FR-059(b)} & Count each eligible current valuation once; show pre-arrival amounts as awaiting receipt without contribution. Keep financial and sync states distinct; infer no payment.\\
\midrule
\rowcolor{TableHeader}
\phantomsection\label{fr:060}\textbf{\identifier{FR-060}} & \textbf{The system shall present authorised prepared indicators offline.}\\*
\phantomsection\label{ac:060:a}\identifier{FR-060(a)} & Show missing/unallocated inputs, available scope, freshness and pending contributions; imply no knowledge of later remote changes.\\
\phantomsection\label{ac:060:b}\identifier{FR-060(b)} & Combine prepared accepted data with local proposals: create once, replace updates and suppress removals. Local arrival enables eligible provisional contributions; only server acceptance changes central totals.\\
\phantomsection\label{ac:060:c}\identifier{FR-060(c)} & Distinguish unavailable data, failures and zero; financial indicators require their grant and prepared data.\\
\midrule
\rowcolor{TableHeader}
\phantomsection\label{fr:061}\textbf{\identifier{FR-061}} & \textbf{The system shall refresh derived totals after source changes without duplicate contributions.}\\*
\phantomsection\label{ac:061:a}\identifier{FR-061(a)} & Recalculate affected origin/destination totals after accepted corrections or late data, including after closure; indicators are read-only.\\
\phantomsection\label{ac:061:b}\identifier{FR-061(b)} & Local proposals affect only labelled provisional results; rejected/unresolved proposals never appear as centrally accepted. Retain unresolved proposals separately.\\
\phantomsection\label{ac:061:c}\identifier{FR-061(c)} & Closure does not freeze results; previous exports remain snapshots.\\
\midrule
\rowcolor{TableHeader}
\phantomsection\label{fr:062}\textbf{\identifier{FR-062}} & \textbf{The system shall export an authorised campaign snapshot as PDF online or offline.}\\*
\phantomsection\label{ac:062:a}\identifier{FR-062(a)} & Include campaign, authorised scope, filters/date basis, generation time, completeness and synchronisation limits.\\
\phantomsection\label{ac:062:b}\identifier{FR-062(b)} & Offline PDFs identify freshness/pending contributions; include money only with permission and prepared data, keeping financial and sync states distinct.\\
\phantomsection\label{ac:062:c}\identifier{FR-062(c)} & Export changes no source records. Provide the generated PDF for authorised download; users may print it through their browser/PDF reader. Detailed Excel export and integrated printing remain possible future enhancements under \decisionref{Q-SCOPE-04}/\decisionref{Q-SCOPE-05}. Worker days/hours are consulted in the profile under \identifier{FR-023}; no dedicated attendance-document export is included.\\
\end{longtable}
}

\Needspace{6\baselineskip}
\subsection{Mapping to Use Cases}
\label{sec:requirements-mapping}
This reverse index records the main and supporting functional requirements
explicitly linked by each use case. Supporting links identify dependencies, not
additional ownership of the same capability.
Each use case also links to its requirements and business rules, while each NFR
identifies the functional capabilities it constrains. These links support impact
analysis and criterion-level verification; they do not establish implementation
or exhaustive test coverage.

{\compacttable
\begin{longtable}{L{0.17}L{0.83}}
\caption{Use-case index to the functional catalogue.}\label{tab:uc-fr-index}\\
\toprule\rowcolor{TableHeader}\textbf{Use case} & \textbf{Main and related functional requirements}\\\midrule\endfirsthead
\toprule\rowcolor{TableHeader}\textbf{Use case} & \textbf{Main and related functional requirements}\\\midrule\endhead
\multicolumn{2}{r}{\footnotesize\itshape Continued on next page}\\\endfoot
\bottomrule\endlastfoot
\hyperref[uc:01]{\identifier{UC-01}} & \hyperref[fr:001]{\identifier{FR-001}}\\
\midrule
\hyperref[uc:02]{\identifier{UC-02}} & \hyperref[fr:002]{\identifier{FR-002}}, \hyperref[fr:003]{\identifier{FR-003}}, \hyperref[fr:004]{\identifier{FR-004}}\\
\midrule
\hyperref[uc:03]{\identifier{UC-03}} & \hyperref[fr:005]{\identifier{FR-005}}, \hyperref[fr:064]{\identifier{FR-064}}\\
\midrule
\hyperref[uc:04]{\identifier{UC-04}} & \hyperref[fr:004]{\identifier{FR-004}}, \hyperref[fr:006]{\identifier{FR-006}}, \hyperref[fr:007]{\identifier{FR-007}}, \hyperref[fr:008]{\identifier{FR-008}}\\
\midrule
\hyperref[uc:05]{\identifier{UC-05}} & \hyperref[fr:009]{\identifier{FR-009}}, \hyperref[fr:010]{\identifier{FR-010}}, \hyperref[fr:011]{\identifier{FR-011}}, \hyperref[fr:012]{\identifier{FR-012}}, \hyperref[fr:013]{\identifier{FR-013}}\\
\midrule
\hyperref[uc:06]{\identifier{UC-06}} & \hyperref[fr:004]{\identifier{FR-004}}, \hyperref[fr:010]{\identifier{FR-010}}, \hyperref[fr:014]{\identifier{FR-014}}, \hyperref[fr:015]{\identifier{FR-015}}, \hyperref[fr:016]{\identifier{FR-016}}, \hyperref[fr:017]{\identifier{FR-017}}, \hyperref[fr:018]{\identifier{FR-018}}, \hyperref[fr:019]{\identifier{FR-019}}\\
\midrule
\hyperref[uc:07]{\identifier{UC-07}} & \hyperref[fr:004]{\identifier{FR-004}}, \hyperref[fr:020]{\identifier{FR-020}}, \hyperref[fr:021]{\identifier{FR-021}}, \hyperref[fr:022]{\identifier{FR-022}}, \hyperref[fr:023]{\identifier{FR-023}}\\
\midrule
\hyperref[uc:08]{\identifier{UC-08}} & \hyperref[fr:024]{\identifier{FR-024}}, \hyperref[fr:025]{\identifier{FR-025}}, \hyperref[fr:026]{\identifier{FR-026}}, \hyperref[fr:027]{\identifier{FR-027}}, \hyperref[fr:029]{\identifier{FR-029}}\\
\midrule
\hyperref[uc:09]{\identifier{UC-09}} & \hyperref[fr:028]{\identifier{FR-028}}\\
\midrule
\hyperref[uc:10]{\identifier{UC-10}} & \hyperref[fr:004]{\identifier{FR-004}}, \hyperref[fr:016]{\identifier{FR-016}}, \hyperref[fr:020]{\identifier{FR-020}}, \hyperref[fr:030]{\identifier{FR-030}}, \hyperref[fr:031]{\identifier{FR-031}}, \hyperref[fr:032]{\identifier{FR-032}}, \hyperref[fr:033]{\identifier{FR-033}}, \hyperref[fr:034]{\identifier{FR-034}}, \hyperref[fr:035]{\identifier{FR-035}}\\
\midrule
\hyperref[uc:11]{\identifier{UC-11}} & \hyperref[fr:033]{\identifier{FR-033}}, \hyperref[fr:036]{\identifier{FR-036}}, \hyperref[fr:037]{\identifier{FR-037}}, \hyperref[fr:061]{\identifier{FR-061}}\\
\midrule
\hyperref[uc:12]{\identifier{UC-12}} & \hyperref[fr:038]{\identifier{FR-038}}, \hyperref[fr:039]{\identifier{FR-039}}, \hyperref[fr:041]{\identifier{FR-041}}, \hyperref[fr:044]{\identifier{FR-044}}\\
\midrule
\hyperref[uc:13]{\identifier{UC-13}} & \hyperref[fr:040]{\identifier{FR-040}}, \hyperref[fr:041]{\identifier{FR-041}}, \hyperref[fr:044]{\identifier{FR-044}}, \hyperref[fr:063]{\identifier{FR-063}}\\
\midrule
\hyperref[uc:14]{\identifier{UC-14}} & \hyperref[fr:020]{\identifier{FR-020}}, \hyperref[fr:041]{\identifier{FR-041}}, \hyperref[fr:042]{\identifier{FR-042}}, \hyperref[fr:043]{\identifier{FR-043}}, \hyperref[fr:044]{\identifier{FR-044}}, \hyperref[fr:045]{\identifier{FR-045}}, \hyperref[fr:063]{\identifier{FR-063}}\\
\midrule
\hyperref[uc:15]{\identifier{UC-15}} & \hyperref[fr:020]{\identifier{FR-020}}, \hyperref[fr:046]{\identifier{FR-046}}, \hyperref[fr:047]{\identifier{FR-047}}, \hyperref[fr:048]{\identifier{FR-048}}, \hyperref[fr:061]{\identifier{FR-061}}\\
\midrule
\hyperref[uc:16]{\identifier{UC-16}} & \hyperref[fr:049]{\identifier{FR-049}}, \hyperref[fr:050]{\identifier{FR-050}}, \hyperref[fr:053]{\identifier{FR-053}}, \hyperref[fr:054]{\identifier{FR-054}}, \hyperref[fr:061]{\identifier{FR-061}}\\
\midrule
\hyperref[uc:17]{\identifier{UC-17}} & \hyperref[fr:051]{\identifier{FR-051}}, \hyperref[fr:052]{\identifier{FR-052}}, \hyperref[fr:054]{\identifier{FR-054}}, \hyperref[fr:058]{\identifier{FR-058}}, \hyperref[fr:061]{\identifier{FR-061}}\\
\midrule
\hyperref[uc:18]{\identifier{UC-18}} & \hyperref[fr:053]{\identifier{FR-053}}, \hyperref[fr:054]{\identifier{FR-054}}, \hyperref[fr:059]{\identifier{FR-059}}, \hyperref[fr:060]{\identifier{FR-060}}, \hyperref[fr:061]{\identifier{FR-061}}, \hyperref[fr:062]{\identifier{FR-062}}\\
\midrule
\hyperref[uc:19]{\identifier{UC-19}} & \hyperref[fr:055]{\identifier{FR-055}}, \hyperref[fr:056]{\identifier{FR-056}}, \hyperref[fr:057]{\identifier{FR-057}}, \hyperref[fr:058]{\identifier{FR-058}}, \hyperref[fr:059]{\identifier{FR-059}}, \hyperref[fr:060]{\identifier{FR-060}}, \hyperref[fr:061]{\identifier{FR-061}}, \hyperref[fr:062]{\identifier{FR-062}}\\
\midrule
\hyperref[uc:20]{\identifier{UC-20}} & \hyperref[fr:038]{\identifier{FR-038}}, \hyperref[fr:040]{\identifier{FR-040}}, \hyperref[fr:041]{\identifier{FR-041}}, \hyperref[fr:042]{\identifier{FR-042}}, \hyperref[fr:043]{\identifier{FR-043}}, \hyperref[fr:044]{\identifier{FR-044}}, \hyperref[fr:045]{\identifier{FR-045}}, \hyperref[fr:063]{\identifier{FR-063}}\\
\end{longtable}}

